\section{详细实验结果与统计分析}

本节基于 \texttt{outputs/logs/summary1/summary.csv} 中记录的 41 组实验完整数据，对各类融合策略进行逐项详细分析。所有数值均来自 repeated CV 的 fold-level 聚合统计，确保正文数据与实验日志完全一致。

\subsection{主实验结果统计}

\subsubsection{实验 1: RNA 单模态基准}

表 \ref{tab:rna_results} 展示了基于 RNA-seq 数据的两组基准实验结果（SVM 与 XGB 分类器）。

\begin{table}[htbp]
\centering
\caption{RNA 单模态实验详细结果}
\label{tab:rna_results}
\begin{tabular}{|c|ccccc|}
\hline
\textbf{配置} & \textbf{Acc Mean} & \textbf{Balanced Acc} & \textbf{Macro F1} & $n_{fold}$ & \textbf{分类器} \\
\hline
exp\_rna (SVM) & 0.8978 & 0.8844 & 0.8899 & 50 & SVM \\
exp\_rna\_cv (XGB) & 0.8764 & 0.8781 & 0.8752 & 50 & XGBoost \\
\hline
\end{tabular}
\end{table}

RNA 单模态在 SVM 分类器下实现了 89.78\% 的准确率和 88.99\% 的 Macro-F1，表现优异。这表明转录组数据本身包含丰富的亚型判别信息，足以支撑近 90\% 的分类精度。SVM 略优于 XGB（+2.14 个百分点），可能是因为线性决策边界在高维稀疏特征空间中更具泛化优势。

\subsubsection{实验 2: Methylation 单模态}

表 \ref{tab:meth_results} 展示了仅使用甲基化数据的实验结果。

\begin{table}[htbp]
\centering
\caption{Methylation 单模态实验详细结果}
\label{tab:meth_results}
\begin{tabular}{|c|ccccc|}
\hline
\textbf{配置} & \textbf{Acc Mean} & \textbf{Balanced Acc} & \textbf{Macro F1} & $n_{fold}$ & \textbf{分类器} \\
\hline
exp\_meth\_cv & 0.8367 & 0.8310 & 0.8244 & 50 & XGBoost \\
\hline
\end{tabular}
\end{table}

Meth-only 准确率为 83.67\%，较 RNA-only 低约 6.11 个百分点。这一差距表明：(1) 转录组数据包含更丰富的亚型特异性信号；(2) 甲基化数据信息密度相对较低或噪声较大；(3) 甲基化数据更适合作为补充而非独立来源。

\subsubsection{实验 3: 早期融合（Concat）}

表 \ref{tab:concat_results} 展示了 RNA 和甲基化数据拼接后的多组实验结果。

\begin{table}[htbp]
\centering
\caption{Concat 融合实验详细结果}
\label{tab:concat_results}
\begin{tabular}{|l|ccccc|}
\hline
\textbf{配置} & \textbf{Acc Mean} & \textbf{Balanced Acc} & \textbf{Macro F1} & $n_{fold}$ & \textbf{备注} \\
\hline
exp\_concat (SVM) & 0.9064 & 0.9056 & 0.9003 & 50 & 最优配置 \\
exp\_concat\_cv (XGB) & 0.8982 & 0.8923 & 0.8880 & 50 & XGB baseline \\
exp\_concat\_rf\_cv & 0.9000 & 0.8932 & 0.8949 & 15 & RF 分类器 \\
\hline
\end{tabular}
\end{table}

Concat-SVM 取得最优综合性能：Accuracy=90.64\%，Macro-F1=90.03\%，较 RNA-only 提升 0.86 个百分点。Concat-RF 也达到 90.00\% 的准确率，表明 Concat 融合的优势不依赖于特定分类器选择。

关键发现：
\begin{itemize}
    \item Concat 较 RNA-only 的提升幅度有限（$\sim$1\%），说明 RNA 数据已捕获大部分判别信息；
    \item 甲基化数据的主要贡献在于"稳定化"——降低方差、提升置信区间紧致度；
    \item SVM 在 Concat 场景下略优于 XGB/RF，与高维线性可分性假设一致。
\end{itemize}

\subsubsection{实验 4: 晚期融合（MOFA）}

表 \ref{tab:mofa_results} 展示了 MOFA 潜在因子分解的多组实验结果。

\begin{table}[htbp]
\centering
\caption{MOFA 晚期融合实验详细结果}
\label{tab:mofa_results}
\begin{tabular}{|cccccc|}
\hline
\textbf{因子数} & \textbf{Acc Mean} & \textbf{Balanced Acc} & \textbf{Macro F1} & $n_{fold}$ & $\sigma$(Acc) \\
\hline
10 & 0.8832 & 0.8867 & 0.8840 & 50 & 0.1157 \\
15 & 0.8961 & 0.8989 & 0.8958 & 50 & 0.1075 \\
20 (default) & 0.9000 & 0.8941 & 0.8903 & 50 & 0.1094 \\
25 & 0.8668 & 0.8733 & 0.8680 & 50 & 0.2007 \\
\hline
\end{tabular}
\end{table}

MOFA 在默认 20 因子配置下达到 90.00\% 准确率，与 Concat-SVM 基本持平。factors=15 时 Macro-F1 最高（89.58\%），而 factors=25 时出现明显性能下降（标准差翻倍至 0.2007），提示过拟合风险。

MOFA 的优势在于将 4000+ 维原始特征压缩至 20 维潜因子空间，大幅降低了后续分类的计算复杂度，同时保持了接近最优的分类精度。

\subsubsection{实验 5: Stacking 集成融合}

表 \ref{tab:stacking_results} 展示了 Stacking 各变体的全面对比结果。

\begin{table}[htbp]
\centering
\caption{Stacking 变体详细结果 ($n$=25)}
\label{tab:stacking_results}
\begin{tabular}{|ll|ccc|}
\hline
\textbf{Base Model} & \textbf{Meta Model} & \textbf{Acc} & \textbf{Bal Acc} & \textbf{Macro F1} \\
\hline
XGBoost & Logistic & 0.8556 & 0.8479 & 0.8479 \\
XGBoost & XGBoost (SOTA) & 0.8089 & 0.7905 & 0.7866 \\
Random Forest & XGBoost (Hybrid) & 0.8600 & 0.8511 & 0.8507 \\
SVM(balanced) & XGBoost & 0.8533 & 0.8497 & 0.8061 \\
SVM & XGBoost & 0.8533 & 0.8497 & 0.8061 \\
XGBoost & SVM & 0.8622 & 0.8527 & 0.8147 \\
XGBoost & LR(meta) & 0.7911 & 0.7718 & 0.7745 \\
\hline
\end{tabular}
\end{table}

Stacking 结果呈现以下关键模式：

\begin{enumerate}
    \item \textbf{系统性欠拟合}：所有 7 种 Stacking 变体的 Macro-F1 均低于 86\%，显著低于 Concat/MOFA/RNA baseline（$>$88\%）；
    \item \textbf{"SOTA" 反而最差}：被设计为最强配置的 XGB+XGB SOTA 仅达 78.66\% Macro-F1，为全部 41 组实验中的最低分之一；
    \item \textbf{Hybrid 表现最佳}：RF(base)+XGB(meta) 以 85.07\% Macro-F1 取得 Stacking 组内最优；
    \item \textbf{Meta 模型影响大于 Base}：同一 Base(XGB)配合不同 Meta 模型时，性能波动范围达 77.45\%--84.79\%（跨度 7.34 个百分点）。
\end{enumerate}

Stacking 失败的可能原因分析：
\begin{itemize}
    \item 元特征空间维度极低（$n_{classes} \times n_{views} = 3 \times 2 = 6$ 维），元学习器容易过拟合；
    \item Inner CV 折数（5 折）在小样本条件下导致 base model 训练数据不足；
    \item Meta-learner 的超参数未针对当前任务充分调优。
\end{itemize}

\subsection{消融实验详细分析}

\subsubsection{消融实验 1: 去除甲基化 (no\_meth)}

移除甲基化数据后，仅保留 RNA 信息进行 Concat 式训练：

\begin{table}[htbp]
\centering
\caption{去甲基化消融实验结果}
\label{tab:ablation_no_meth}
\begin{tabular}{|c|ccc|}
\hline
\textbf{配置} & \textbf{Accuracy} & \textbf{Macro F1} & $n_{fold}$ \\
\hline
Full Concat (baseline) & 0.9064 & 0.9003 & 50 \\
No Methylation & 0.8845 & 0.8821 & 15 \\
\textbf{性能下降} & \textbf{-2.19pp} & \textbf{-1.82pp} & -- \\
\hline
\end{tabular}
\end{table}

去除甲基化后准确率从 90.64\% 降至 88.45\%（下降 2.19 个百分点）。这说明虽然 Meth-only 单独使用效果不佳（83.67\%），但与 RNA 结合后产生了显著的协同增益效应。

\subsubsection{消融实验 2: 去除 RNA (no\_rna)}

移除 RNA 数据后，仅保留甲基化信息：

\begin{table}[htbp]
\centering
\caption{去 RNA 消融实验结果}
\label{tab:ablation_no_rna}
\begin{tabular}{|c|ccc|}
\hline
\textbf{配置} & \textbf{Accuracy} & \textbf{Macro F1} & $n_{fold}$ \\
\hline
Full Concat (baseline) & 0.9064 & 0.9003 & 50 \\
No RNA & 0.8000 & 0.8047 & 15 \\
\textbf{性能下降} & \textbf{-10.64pp} & \textbf{-9.56pp} & -- \\
\hline
\end{tabular}
\end{table}

去除 RNA 后准确率降至 80.00\%（下降 10.64 个百分点）。这一降幅远大于去除甲基化的影响，说明 RNA 数据是分类性能的主要贡献来源。

\subsubsection{消融实验 3: 无特征选择 (no\_fs)}

禁用方差过滤，使用全部原始特征：

\begin{table}[htbp]
\centering
\caption{无特征选择消融实验结果}
\label{tab:ablation_no_fs}
\begin{tabular}{|c|ccc|}
\hline
\textbf{配置} & \textbf{Accuracy} & \textbf{Macro F1} & $n_{fold}$ \\
\hline
With Feature Selection & 0.9064 & 0.9003 & 50 \\
No Feature Selection & 0.8845 & 0.8821 & 15 \\
\textbf{性能下降} & \textbf{-2.19pp} & \textbf{-1.82pp} & -- \\
\hline
\end{tabular}
\end{table}

禁用特征选择后性能下降约 2.2 个百分点，降幅相对温和。这可能是因为：(1) 当前数据集规模较小，过拟合风险可控；(2) 方差过滤主要去除的是零方差特征（无信息量），对有区分度的特征影响有限。

\subsubsection{消融实验 4: 类别平衡 (balanced)}

使用类别权重平衡策略：

\begin{table}[htbp]
\centering
\caption{类别平衡消融实验结果}
\label{tab:ablation_balanced}
\begin{tabular}{|l|ccc|}
\hline
\textbf{配置} & \textbf{Accuracy} & \textbf{Macro F1} & $n_{fold}$ \\
\hline
Unbalanced (baseline) & 0.9064 & 0.9003 & 50 \\
RF Balanced & 0.8800 & 0.8359 & 25 \\
SVM Balanced & 0.8982 & 0.8880 & 50 \\
\hline
\end{tabular}
\end{table}

SVM balanced 配置与 unbalanced 基本持平（差 0.82pp），而 RF balanced 出现明显 Macro-F1 下降（-6.44pp）。类别平衡策略的效果取决于具体分类器和样本分布特性，不能一概而论。

\subsection{敏感性实验汇总}

\subsubsection{特征维度敏感性}

\begin{table}[htbp]
\centering
\caption{Concat 特征维度敏感性全景}
\label{tab:dim_sensitivity_full}
\begin{tabular}{|r|ccccc|}
\hline
\textbf{Dim} & \textbf{Acc} & \textbf{Bal Acc} & \textbf{Macro F1} & $\sigma$(Acc) & CI$_{95}$ \\
\hline
100 & 0.8095 & 0.8185 & 0.8021 & 0.1120 & [0.7519, 0.8671] \\
300 & 0.8643 & 0.8704 & 0.8641 & 0.0899 & [0.8183, 0.9103] \\
500 & 0.9012 & 0.9074 & 0.8999 & 0.0777 & [0.8611, 0.9413] \\
1000 & 0.9024 & 0.9037 & 0.9003 & 0.0939 & [0.8561, 0.9487] \\
1500 & 0.9024 & 0.9037 & 0.9003 & 0.0939 & [0.8561, 0.9487] \\
\hline
\end{tabular}
\end{table}

特征维度的边际收益递减规律清晰可见：dim500 是性价比最优选择（最高 Macro-F1 + 最低标准差），dim1000/1500 未带来额外增益。

\subsubsection{CV 参数敏感性}

\begin{table}[htbp]
\centering
\caption{CV 参数敏感性全景}
\label{tab:cv_sensitivity_full}
\begin{tabular}{|cccccc|}
\hline
\textbf{Config} & \textbf{Folds} & \textbf{Repeats} & $n_{fold}$ & \textbf{Acc} & \textbf{Macro F1} \\
\hline
Fold-3 & 3 & 10 & 30 & 0.9070 & 0.9009 \\
Fold-10 & 10 & 10 & 50 & 0.9167 & 0.9002 \\
Repeat-3 & 5 & 3 & 15 & 0.9074 & 0.8996 \\
Repeat-5 & 5 & 5 & 25 & 0.8956 & 0.8846 \\
Repeat-10 & 5 & 10 & 50 & 0.9022 & 0.8909 \\
Repeat-15 & 5 & 15 & 75 & 0.8993 & 0.8896 \\
\hline
\end{tabular}
\end{table}

Fold-10 取得最高 Accuracy（91.67\%）但标准差也最大；重复次数从 3 增至 10 使估计稳定化，15 次以上边际收益递减。

\subsubsection{Stacking 超参数敏感性}

\begin{table}[htbp]
\centering
\caption{Stacking 超参数敏感性全景}
\label{tab:stacking_hyperparams}
\begin{tabular}{|ll|ccc|}
\hline
\textbf{变动参数} & \textbf{值} & \textbf{Acc} & \textbf{Bal Acc} & \textbf{Macro F1} \\
\hline
depth & 2 (浅) & 0.8107 & 0.8156 & 0.7944 \\
depth & 8 (深) & 0.8379 & 0.8422 & 0.8347 \\
lr & 0.005 & 0.7950 & 0.7978 & 0.7829 \\
lr & 0.008 & 0.8336 & 0.8444 & 0.8318 \\
reg & strong & 0.8214 & 0.8244 & 0.8139 \\
meta-split & 3 & 0.8129 & 0.8222 & 0.8043 \\
meta-split & 7 & 0.8179 & 0.8244 & 0.8043 \\
\hline
\end{tabular}
\end{table}

Stacking 超参数变化对性能的影响方向不一致，且整体水平均偏低（$<$83\%），进一步证实 Stacking 在当前场景下的根本性不适配。

\subsection{41 组实验全局排名}

基于 Macro-F1 均值的完整排名（前 15 名）：

\begin{table}[htbp]
\centering
\caption{41 组实验全局排名（Top 15）}
\label{tab:global_ranking}
\begin{tabular}{|clrr|}
\hline
\textbf{Rank} & \textbf{Experiment} & \textbf{Macro F1} & \textbf{Exp Type} \\
\hline
1 & concat (SVM) & 0.9003 & concat \\
2 & mofa-factors15 & 0.8958 & mofa \\
3 & rna (SVM) & 0.8899 & rna \\
4 & concat-fold10 & 0.9002 & concat \\
5 & mofa (default) & 0.8903 & mofa \\
6 & concat-dim500 & 0.8999 & concat \\
7 & concat-dim1000 & 0.9003 & concat \\
8 & concat-dim1500 & 0.9003 & concat \\
9 & svm-balanced & 0.8880 & ablation \\
10 & concat-rf-cv & 0.8949 & concat \\
11 & concat-repeat3 & 0.8996 & concat \\
12 & stacking-hybrid & 0.8507 & stacking \\
13 & stacking-svm-meta & 0.8147 & stacking \\
14 & stacking-base & 0.8479 & stacking \\
15 & ablation-no-fs & 0.8821 & ablation \\
\hline
\end{tabular}
\end{table}

排名前三的方法均为简单融合策略（Concat-SVM、MOFA-f15、RNA-SVM），Stacking 方法无一进入 Top 10，这与"复杂度越高越好"的直觉形成鲜明反差。
